% BHU PYQ -- Manually transcribed & error-corrected
% Paper: BCH-315 : Auditing
% Exam: B.Com. (Hons.) Semester V Examination, 2022-23

\documentclass[11pt,a4paper]{article}
\usepackage[a4paper,top=2.5cm,bottom=2.5cm,left=2.8cm,right=2.8cm,headheight=14pt]{geometry}
\usepackage{amsmath,amssymb}
\usepackage[T1]{fontenc}
\usepackage[utf8]{inputenc}
\usepackage{lmodern,microtype,fancyhdr,enumitem,hyperref,lastpage,parskip}

\hypersetup{
    colorlinks=true,
    urlcolor=black,
    linkcolor=black,
    pdftitle={BCH-315: Auditing -- BHU B.Com. (Hons.) Sem V 2022-23}
}

\pagestyle{fancy}
\fancyhf{}
\renewcommand{\headrulewidth}{0.4pt}
\renewcommand{\footrulewidth}{0.4pt}
\fancyhead[L]{\small   \textit{Banaras Hindu University}}
\fancyhead[R]{\small   \textit{BCH-315: Auditing}}
\fancyfoot[C]{\small Page  \thepage\ of \pageref{LastPage}}


\newlist{parts}{enumerate}{2}
\setlist[parts,1]{label=(\alph*),leftmargin=*,itemsep=5pt,topsep=3pt}
\setlist[parts,2]{label=(\roman*),leftmargin=*,itemsep=3pt,topsep=2pt}

\newcommand{\pts}[1]{\hfill   \textit{[#1]}}


\begin{document}


\begin{center}
  {\large\bfseries BANARAS HINDU UNIVERSITY}\\[2pt]
  {
\normalsize B.Com. (Hons.) --- Semester V Examination, 2022-23}\\[6pt]
  \rule{0.8\linewidth}{0.5pt}\\[6pt]
  {\large\bfseries Paper No. BCH-315: Auditing}\\[4pt]
  {
\normalsize Time: 3 Hours\hfill Maximum Marks: 70}
\end{center}

\rule{\linewidth}{0.4pt}

\smallskip



\noindent   \textit{Instructions: Attempt all questions. All questions carry equal marks (14 marks each).}

\bigskip




\noindent   \textbf{Question 1.}
What do you mean by Auditing? Discuss its objectives.\pts{14}

\centerline{   \textbf{OR}}

Write a note on preparation before commencement of audit.\pts{14}

\medskip\rule{\linewidth}{0.2pt}




\noindent   \textbf{Question 2.}
Explain the various points that should be taken into consideration while adopting techniques of vouching.\pts{14}

\centerline{   \textbf{OR}}

``Vouching is the essence of auditing.'' Comment on this statement and explain the points to be borne in mind while examining the vouchers.\pts{14}

\medskip\rule{\linewidth}{0.2pt}




\noindent   \textbf{Question 3.}
What do you understand by verification? Explain the objectives of verification.\pts{14}

\centerline{   \textbf{OR}}

Distinguish between the terms valuation and verification. ``An auditor is not a valuer.'' Discuss.\pts{14}

\medskip\rule{\linewidth}{0.2pt}




\noindent   \textbf{Question 4.}
Describe the provisions of the Companies Act regarding qualifications, disqualifications and removal of auditors.\pts{14}

\centerline{   \textbf{OR}}

Point out the powers and duties of a Company Auditor.\pts{14}

\medskip\rule{\linewidth}{0.2pt}




\noindent   \textbf{Question 5.}
Distinguish between clean and qualified report. Give the specimen of an imaginary clean report.\pts{14}

\centerline{   \textbf{OR}}

What are the contents of an Auditor's Report? Discuss.\pts{14}

\vfill
\rule{\linewidth}{0.4pt}

\begin{center}\small
\href{https://bhu-exam.vercel.app/}{Click here to visit our website}\\[4pt]
\href{https://me-aryan.vercel.app/}{Click here to visit our contributor}\\[4pt]
\href{https://quashan-me.vercel.app/}{Click here to visit our contributor}
\end{center}

\end{document}
